\documentclass[11pt]{article}
\usepackage[T1]{fontenc}
\usepackage[utf8]{inputenc}
\usepackage{lmodern}
\usepackage{microtype}
\usepackage[margin=0.92in]{geometry}
\usepackage{amsmath,amssymb,amsthm,mathtools}
\usepackage{booktabs,array,longtable,tabularx}
\usepackage{enumitem}
\usepackage{xcolor}
\usepackage{hyperref}
\usepackage{listings}
\usepackage{fancyhdr}
\usepackage{enumitem}
\usepackage{bm}
\usepackage{needspace}
\usepackage{url}

\hypersetup{
 colorlinks=true,
 linkcolor=blue!55!black,
 citecolor=blue!55!black,
 urlcolor=blue!55!black,
 pdftitle={Von Mangoldt and Euler totient identities in intuitionistic primitive recursive arithmetic},
 pdfauthor={OpenAI}
}
\pagestyle{fancy}
\setlength{\headheight}{14pt}
\fancyhf{}
\lhead{Mangoldt and totient in iPRA}
\rhead{Constructive proof development}
\cfoot{\thepage}

\newtheorem{theorem}{Theorem}[section]
\newtheorem{lemma}[theorem]{Lemma}
\newtheorem{corollary}[theorem]{Corollary}
\newtheorem{proposition}[theorem]{Proposition}
\newtheorem{claim}[theorem]{Claim}
\theoremstyle{definition}
\newtheorem{definition}[theorem]{Definition}
\newtheorem{remark}[theorem]{Remark}
\newtheorem{convention}[theorem]{Convention}
\newtheorem{example}[theorem]{Example}

\newcommand{\N}{\mathbb N}
\newcommand{\Z}{\mathbb Z}
\newcommand{\Q}{\mathbb Q}
\newcommand{\ipra}{\mathrm{iPRA}}
\newcommand{\ph}{\varphi}
\newcommand{\rad}{\operatorname{rad}}
\newcommand{\lcmu}{\operatorname{lcmUpto}}
\newcommand{\ord}{\operatorname{ord}}
\newcommand{\Prime}{\operatorname{Prime}}
\newcommand{\PP}{\operatorname{PrimePow}}
\newcommand{\wtlog}{L_w}
\newcommand{\wLambda}{\Lambda_w}
\newcommand{\wPsi}{\Psi_w}
\newcommand{\eps}{\varepsilon}
\newcommand{\code}[1]{\nolinkurl{#1}}
\newcommand{\StatusI}{\textsf{I}}
\newcommand{\StatusF}{\textsf{F}}
\newcommand{\StatusA}{\textsf{A}}
\newcommand{\StatusO}{\textsf{O}}
\newcommand{\ind}[1]{\mathbf 1_{#1}}

\setlist[itemize]{leftmargin=1.5em,itemsep=0.25em,topsep=0.35em}
\setlist[enumerate]{leftmargin=1.8em,itemsep=0.25em,topsep=0.35em}
\setlength{\parskip}{0.25em}
\setlength{\parindent}{1.2em}
\setlength{\emergencystretch}{2em}
\Urlmuskip=0mu plus 1mu

\lstdefinelanguage{Lean}{
 morekeywords={def,theorem,lemma,namespace,end,by,have,show,from,import,open,section,variable,where,if,then,else,let,in,match,with,case,induction,exact,refine,simp,rw,constructor,structure,inductive,axiom,noncomputable},
 sensitive=true,
 morecomment=[l]{--},
 morecomment=[s]{/-}{-/},
 morestring=[b]"
}
\lstset{language=Lean,basicstyle=\ttfamily\footnotesize,columns=fullflexible,breaklines=true,frame=single,showstringspaces=false,backgroundcolor=\color{black!2},xleftmargin=0.3em,xrightmargin=0.3em}

\title{\textbf{The von Mangoldt Function and Euler's Totient Function}\\
\large Finite, constructive proofs in intuitionistic primitive recursive arithmetic}
\author{Research note prepared for Vladimir Reshetnikov}
\date{15 November 2026}

\begin{document}
\maketitle

\begin{abstract}
This note reconstructs, in a form suitable for intuitionistic primitive
recursive arithmetic (\(\ipra\)), the arithmetic content of the von Mangoldt
and Euler totient reference pages listed in the bibliography.  The central
point is that the elementary theorems do not need classical logic, unbounded
choice, completed infinite sets, or real analysis.  Divisibility, primality,
valuation, finite factorization, the M\"obius function, bounded sums and
products, reduced residue systems, and finite Dirichlet convolution are all
primitive recursive.  Their proofs can therefore be expressed by
quantifier-free induction.

For the von Mangoldt function we separate its exact arithmetic skeleton from
the real logarithm.  Given any primitive-recursive weight \(w(p)\) on primes,
we define the completely additive weight
\(L_w(n)=\sum_p v_p(n)w(p)\) and the weighted Mangoldt function
\(\Lambda_w(p^a)=w(p)\).  All divisor sums, M\"obius inversion formulas,
least-common-multiple identities, Chebyshev sums, and formal Dirichlet-series
identities are then literal \(\ipra\) theorems.  Taking \(w(p)=\log p\) is a
constructive-analysis specialization, implemented internally at any requested
finite precision by primitive-recursive rational bounds.

For Euler's totient we prove the counting definition, prime and prime-power
formulae, multiplicativity and the product formula, divisor-sum and M\"obius
identities, parity and divisibility results, power/gcd/lcm/radical identities,
Euler's theorem, RSA correctness, Menon's identity, the root-of-unity identity,
Lambert and Dirichlet coefficient identities, exact summatory formulae,
several elementary bounds, and finite witness forms of the limsup/liminf
statements.  The archive also supplies a theorem-by-theorem status ledger, a
Lean formalization blueprint tied to the current \code{ConciseIPRA}
infrastructure, host-level Lean scaffolding, and an executable finite checker.
No claim is made that the supplied Lean files are a completed formalization.
\end{abstract}

\tableofcontents
\newpage

\section{Scope, status labels, and the constructive reading}

\subsection{What it means here to prove a reference-page statement in \(\ipra\)}

The language of ordinary \(\ipra\) has one natural-number sort, a symbol for
each primitive-recursive function, equality, intuitionistic first-order logic,
and induction for quantifier-free formulas.  Consequently a theorem is in the
native scope of \(\ipra\) when it can be written using natural-number codes,
primitive-recursive operations, and quantifier-free induction.  Finite sets
are represented by bounded characteristic functions; their cardinalities,
sums, and products are bounded folds.  Integers are represented by pairs of
naturals, and rational equalities or inequalities are denominator-cleared.

Some reference-page formulas use real logarithms, complex roots of unity,
infinite series, limits, or asymptotic notation.  These are not literally
formulas of first-order \(\ipra\).  We therefore distinguish four statuses.

\begin{center}
\begin{tabularx}{0.94\textwidth}{>{\bfseries}c X}
\toprule
Label & Meaning \\
\midrule
\StatusI & A literal, denominator-cleared natural/integer statement whose proof uses primitive-recursive constructions and quantifier-free induction. \\
\StatusF & An exact finite coefficient, polynomial-remainder, or convolution statement in \(\ipra\) which is the algebraic content of an infinite analytic notation. \\
\StatusA & A constructive-analysis corollary obtained uniformly from finite \(\ipra\) certificates, such as rational upper and lower approximations to logarithms. \\
\StatusO & Not claimed here: genuinely infinitary/complex-analytic material, an open problem, or a result whose complete constructive reduction would require a separate large development. \\
\bottomrule
\end{tabularx}
\end{center}

This distinction is mathematical rather than cosmetic.  For example,
\(\sum_{d\mid n}\Lambda(d)=\log n\) contains a real logarithm.  Its exact
arithmetic content is the stronger weighted identity
\[
 \sum_{d\mid n}\Lambda_w(d)=L_w(n)
\]
for every primitive-recursive prime weight \(w\).  The usual formula follows
by interpreting \(w(p)\) as \(\log p\).  The weighted identity is status
\StatusI; the real specialization is status \StatusA.

\subsection{Coverage summary}

The report proves or reduces to exact finite form the following principal
families.

\begin{longtable}{@{}p{0.07\textwidth}p{0.78\textwidth}p{0.07\textwidth}@{}}
\toprule
No. & Statement family & Status \\
\midrule
\endfirsthead
\toprule No. & Statement family & Status \\
\midrule
\endhead
T1 & Definition of \(\varphi(n)\) as the number of residues below \(n\) coprime to \(n\); primitive recursiveness; \(\varphi(0)=0,\varphi(1)=1\). & \StatusI\\
T2 & \(\varphi(p)=p-1\), \(\varphi(p^a)=p^{a-1}(p-1)\), multiplicativity, and Euler's product formula. & \StatusI\\
T3 & \(\sum_{d\mid n}\varphi(d)=n\), M\"obius inversion, gcd-fiber counting, and summatory floor-square identities. & \StatusI\\
T4 & Power, divisibility, gcd/lcm, doubling, radical, and squarefree-divisor identities. & \StatusI\\
T5 & Evenness, powers of two dividing \(\varphi(n)\), sum of totatives, and Menon's identity. & \StatusI\\
T6 & Euler's theorem, Fermat's theorem, RSA correctness, and \(m\mid\varphi(a^m-1)\). & \StatusI\\
T7 & Root-of-unity/cosine identity through an exact cyclotomic polynomial congruence. & \StatusF\\
T8 & Lambert-series and Dirichlet-series formulas through exact coefficient and convolution identities. & \StatusF\\
T9 & Bounds \(\varphi(n)^2\ge n/2\), the sharper square-root exception theorem, composite upper bound, and \(\varphi(n)\sigma(n)/n^2>1/2\). & \StatusI\\
T10 & Finite witness forms of \(\limsup\varphi(n)/n=1\) and \(\liminf\varphi(n)/n=0\). & \StatusI/\StatusA\\
T11 & Odd nontotients, \(\varphi(n)=2\) classification, and the arithmetic restriction on prime-power constructibility. & \StatusI\\
M1 & Primitive-recursive recognition of prime powers and the weighted Mangoldt function \(\Lambda_w\). & \StatusI\\
M2 & Divisor sum and all standard M\"obius inversion/convolution identities. & \StatusI\\
M3 & \(\exp(\Lambda(n))\) as an lcm quotient, in the exact weighted prime-factor form. & \StatusI/\StatusA\\
M4 & Chebyshev sum \(\Psi_w(N)=L_w(\operatorname{lcm}(1,\dots,N))\). & \StatusI\\
M5 & Formal logarithmic derivative of the zeta Dirichlet series and the prime-power coefficients of \(\log\zeta\). & \StatusF\\
M6 & Finite Abel summation and explicit elementary linear bounds for the Chebyshev function. & \StatusI/\StatusA\\
M7 & \(\sum_{n\le N}\Lambda(n)/n=\log N+O(1)\) with explicit constants, via denominator-cleared finite inequalities. & \StatusI/\StatusA\\
M8 & Generalized Mangoldt functions \(\Lambda_k\), including \(\Lambda_1=\Lambda\). & \StatusI/\StatusF\\
M9 & Riesz means, explicit formulas involving zeta zeros, RH equivalences, and oscillation theorems. & \StatusO\\
\bottomrule
\end{longtable}

\section{Primitive-recursive arithmetic toolkit}

\subsection{Bounded operations}

\begin{lemma}[Bounded-fold closure]\label{lem:bounded-fold}
Let \(f(\vec x,t)\) be primitive recursive.  Then the functions
\[
 \sum_{t<b}f(\vec x,t),\qquad
 \prod_{t<b}f(\vec x,t),\qquad
 \min\{t<b:P(\vec x,t)\},\qquad
 \#\{t<b:P(\vec x,t)\}
\]
are primitive recursive whenever \(P\) has a primitive-recursive
characteristic function.  The minimum is totalized by returning \(b\) when
there is no witness.
\end{lemma}

\begin{proof}
The sum and product are ordinary primitive recursions on \(b\).  A bounded
minimum carries a register initially equal to \(b\); at stage \(t\), it is
replaced by \(t\) exactly when no earlier witness has been stored and the
Boolean characteristic of \(P(\vec x,t)\) is one.  A cardinality is the sum of
the characteristic function.  Every update is obtained from composition,
primitive recursion, truncated subtraction, and zero tests.  The defining
invariants are equations between primitive-recursive terms, so their
correctness proofs use quantifier-free induction.  No instance of excluded
middle beyond decidable equality/order on naturals occurs.
\end{proof}

Division and remainder can be obtained by a bounded scan for the largest
\(q\le x\) satisfying \(qy\le x\), with the convention \(x/0=0\).  Thus
\(d\mid n\), gcd, lcm, primality, bounded lists of divisors, and exact quotient
are primitive recursive.  Gcd may also be implemented by Euclid's algorithm
with a fuel counter \(x+y+1\); its decreasing invariant is quantifier-free.

\subsection{Valuations and a canonical finite factorization}

For \(p,n\in\N\), put
\[
 v_p(n)=\max\{e\le n:p^e\mid n\},
\]
with \(v_p(0)=0\) for totality.  Let \(\ell(n)\) be the least divisor of
\(n>1\) in \([2,n]\), and set \(\ell(0)=\ell(1)=1\).

\begin{lemma}[Least-factor trace]\label{lem:factor-trace}
There is a primitive-recursive trace \(x_0(n),\ldots,x_n(n)\) such that
\[
 x_0=n,
 \qquad
 x_{i+1}=\begin{cases}x_i/\ell(x_i),&x_i>1,\\1,&x_i\le1,
 \end{cases}
\]
and, for \(n>0\), the trace reaches \(1\) by stage \(n\).  Every nontrivial
factor selected by the trace is prime, and
\[
 n=\prod_{p=2}^{n}p^{v_p(n)}.
\]
All assertions admit quantifier-free induction proofs in \(\ipra\).
\end{lemma}

\begin{proof}
If \(x_i>1\), then \(2\le\ell(x_i)\le x_i\), so
\(x_{i+1}<x_i\).  The invariant
\[
 x_i+i\le n+1\quad\text{whenever }x_i>1
\]
is immediate by induction and implies termination within \(n\) stages.  If
\(\ell(x_i)=ab\) with \(1<a<\ell(x_i)\), then \(a\mid x_i\), contradicting
minimality; hence \(\ell(x_i)\) is prime.  Multiplying the exact division
equations along the trace reconstructs \(n\).  Counting occurrences of each
prime gives the valuation form.  Every counter, product, and assertion is
bounded by a primitive-recursive term.
\end{proof}

The usual uniqueness of factorization can be proved without a classical
minimal-counterexample argument.  If
\(\prod p^{a_p}=\prod p^{b_p}\), apply the primitive-recursive valuation
\(v_q\) to both sides; the elementary rules
\(v_q(xy)=v_q(x)+v_q(y)\) and \(v_q(p^a)=a\ind{p=q}\) give \(a_q=b_q\).
These valuation rules are proved by bounded induction on the exponents.

\subsection{M\"obius function and finite inversion}

Encode signed integers as pairs \((u,v)\), representing \(u-v\).  Define
\(\mu(1)=1\).  For \(n>1\), let \(\mu(n)=0\) if some prime square divides
\(n\), and otherwise let \(\mu(n)=(-1)^r\), where \(r\) is the number of
prime divisors.  This is primitive recursive.

\begin{lemma}[M\"obius cancellation]\label{lem:mobius-cancel}
For every positive \(n\),
\[
 \sum_{d\mid n}\mu(d)=\begin{cases}1,&n=1,\\0,&n>1.
 \end{cases}
\]
Consequently, for natural- or integer-valued primitive-recursive functions
\(f,g\),
\[
 g(n)=\sum_{d\mid n}f(d)
 \quad\Longleftrightarrow\quad
 f(n)=\sum_{d\mid n}\mu(d)g(n/d).
\]
\end{lemma}

\begin{proof}
If \(n=\prod_{i<r}p_i^{a_i}\), only squarefree divisors contribute.  Choosing
a squarefree divisor is choosing a bit vector \(e\in\{0,1\}^r\), and its
contribution is \((-1)^{\sum e_i}\).  The sum is
\(\prod_{i<r}(1-1)\), equal to one when \(r=0\) and zero otherwise.  This
finite distributive identity is proved by induction on \(r\).

For inversion, substitute the proposed expression and swap two bounded sums:
\[
 \sum_{d\mid n}\mu(d)\sum_{e\mid n/d}f(e)
 =\sum_{e\mid n}f(e)\sum_{d\mid n/e}\mu(d).
\]
The inner sum is one exactly when \(e=n\), so only \(f(n)\) remains.  The
reverse implication is identical with the roles exchanged.  Sum swapping is a
bounded-rectangle induction, not an appeal to an infinite Fubini theorem.
\end{proof}

\subsection{Finite logarithms and prime weights}

The real number \(\log p\) is not a natural number and cannot be a primitive
symbol of first-order \(\ipra\).  The algebra needs only complete additivity,
so we isolate it.

\begin{definition}[Prime weight and finite logarithm]\label{def:weight}
Let \(w:\N\to\N\) (or a signed/rational code) be primitive recursive.  Define
\[
 L_w(0)=0,
 \qquad
 L_w(n)=\sum_{p=2}^{n}\ind{\Prime(p)}v_p(n)w(p)
 \quad(n>0).
\]
Then \(L_w(1)=0\) and \(L_w(ab)=L_w(a)+L_w(b)\) for positive \(a,b\).
\end{definition}

\begin{proof}
The definition is a bounded sum.  Complete additivity follows termwise from
\(v_p(ab)=v_p(a)+v_p(b)\), followed by finite distributivity.  This proof is
uniform in the code of \(w\).
\end{proof}

For the classical specialization one may use rational interval codes
\(\operatorname{LogLo}(p,t)\le\log p\le\operatorname{LogHi}(p,t)\) of width
at most \(2^{-t}\).  A convenient construction reduces \(p\) by powers of two
and uses
\[
 \log x=2\sum_{j=0}^{M-1}\frac{z^{2j+1}}{2j+1}+R_M,
 \qquad z=\frac{x-1}{x+1}.
\]
A finite rational sum with a geometric bound on \(R_M\).  Numerators,
denominators, powers, and the required \(M\) are primitive recursive.
Therefore every finite identity in \(L_w\) yields certified lower and upper
bounds for its logarithmic interpretation at arbitrary precision.

\section{Euler's totient: definition and factorization}

\begin{definition}[Totient]\label{def:totient}
For \(n\in\N\), define
\[
 \varphi(n)=\sum_{a<n}\ind{\gcd(a,n)=1}.
\]
We use the totalization \(\varphi(0)=0\).  Thus \(\varphi(1)=1\), because the
only \(a<1\) is \(0\), and \(\gcd(0,1)=1\).
\end{definition}

The definition is primitive recursive by Lemma~\ref{lem:bounded-fold}.  It is
extensionally identical to the usual number of positive integers at most
\(n\) that are relatively prime to \(n\), since for \(n>1\), zero is not
coprime to \(n\), while the endpoint \(n\) is not coprime either.

\begin{remark}[The value at zero]
The requested sources concern positive inputs.  Extensions at zero are conventional: some tables set \(\varphi(0)=1\), while the Wolfram Language and the direct bounded-count definition use zero.  We adopt \(\varphi(0)=0\), matching \code{Nat.totient} and avoiding any effect on the positive-integer theorems.
\end{remark}

\begin{theorem}[Prime and prime-power formulae]\label{thm:phi-prime-power}
If \(p\) is prime and \(a\ge1\), then
\[
 \varphi(p)=p-1,
 \qquad
 \varphi(p^a)=p^a-p^{a-1}=p^{a-1}(p-1).
\]
\end{theorem}

\begin{proof}
Every \(1\le x<p\) is coprime to the prime \(p\), so exactly the \(p-1\)
nonzero residues are counted.  For \(p^a\), a residue fails to be coprime
exactly when it is divisible by \(p\).  The map
\(y\mapsto py\) is a bijection from \([0,p^{a-1})\) to the multiples of
\(p\) below \(p^a\).  Hence there are \(p^{a-1}\) excluded residues and
\(p^a-p^{a-1}\) included residues.  Injectivity, range, and the cardinality
calculation are bounded equations and inequalities.
\end{proof}

\begin{theorem}[Constructive Chinese remainder theorem]\label{thm:crt}
If \(m,n>0\) and \(\gcd(m,n)=1\), the map
\[
 x\bmod mn\longmapsto(x\bmod m,x\bmod n)
\]
is a primitive-recursive bijection between residue codes
\([0,mn)\) and \([0,m)\times[0,n)\).  It restricts to a bijection of reduced
residue systems.
\end{theorem}

\begin{proof}
The extended Euclidean algorithm with bounded fuel computes \(u,v\in\Z\)
with \(um+vn=1\).  Given \(a<m,b<n\), put
\[
 x\equiv b(um)+a(vn)\pmod{mn}.
\]
After taking the canonical remainder.  Modulo \(m\), the first term is zero
and \(vn\equiv1\), so \(x\equiv a\); modulo \(n\), similarly
\(x\equiv b\).  If two numbers in \([0,mn)\) have both residues equal, their
difference is divisible by \(m\) and \(n\); coprimality gives divisibility by
\(mn\), and the size bound gives equality.  Finally,
\(\gcd(x,mn)=1\) iff \(\gcd(x,m)=\gcd(x,n)=1\), proved by the elementary
prime-divisor criterion.  All searches and canonical remainders are primitive
recursive.
\end{proof}

\begin{theorem}[Multiplicativity and Euler product]\label{thm:phi-product}
If \(\gcd(m,n)=1\), then
\[
 \varphi(mn)=\varphi(m)\varphi(n).
\]
For every positive \(n=\prod_{p\mid n}p^{a_p}\),
\[
 \boxed{\displaystyle
 \varphi(n)=\prod_{p\mid n}p^{a_p-1}(p-1)
 =n\prod_{p\mid n}\left(1-\frac1p\right).}
\]
The second display is interpreted internally by its denominator-cleared first
product.
\end{theorem}

\begin{proof}
The restricted CRT bijection has domain size \(\varphi(mn)\) and codomain size
\(\varphi(m)\varphi(n)\), proving multiplicativity.  Distinct prime powers in
the factorization of \(n\) are pairwise coprime, so finite induction over the
factor trace and Theorem~\ref{thm:phi-prime-power} yield the first product.
Multiplying and dividing each local factor by \(p^{a_p}\) gives the familiar
rational notation; no division is needed in the internal theorem.
\end{proof}

\begin{example}
The formula gives
\[
 \varphi(24)=\varphi(2^3)\varphi(3)=4\cdot2=8,
\]
corresponding to \(1,5,7,11,13,17,19,23\).
\end{example}

\section{Divisor sums, M\"obius inversion, and exact summatory formulae}

\begin{theorem}[Gcd fibers]\label{thm:gcd-fiber}
If \(d\mid n\), then
\[
 \#\{0\le k<n:\gcd(k,n)=d\}=\varphi(n/d).
\]
\end{theorem}

\begin{proof}
Write \(n=dm\).  Every \(k<n\) with gcd \(d\) is uniquely \(k=da\) with
\(a<m\), and
\(\gcd(da,dm)=d\gcd(a,m)\).  Thus the gcd is \(d\) exactly when
\(\gcd(a,m)=1\).  The maps \(a\mapsto da\) and exact division by \(d\) are
inverse bounded maps.
\end{proof}

\begin{theorem}[Divisor-sum identity]\label{thm:sum-phi}
For every positive \(n\),
\[
 \boxed{\displaystyle\sum_{d\mid n}\varphi(d)=n.}
\]
\end{theorem}

\begin{proof}
Partition the \(n\) residues \(k<n\) by \(d=\gcd(k,n)\).  By
Theorem~\ref{thm:gcd-fiber}, the fiber at \(d\) has size \(\varphi(n/d)\).
Hence
\[
 n=\sum_{d\mid n}\varphi(n/d)=\sum_{e\mid n}\varphi(e),
\]
where the last equality uses the primitive-recursive involution
\(d\mapsto n/d\) on the divisor list.
\end{proof}

\begin{corollary}[M\"obius formula]\label{cor:phi-mobius}
For positive \(n\),
\[
 \boxed{\displaystyle
 \varphi(n)=\sum_{d\mid n}\mu(d)\frac nd
 =\sum_{d\mid n}d\,\mu(n/d).}
\]
Every quotient displayed is exact.
\end{corollary}

\begin{proof}
Apply Lemma~\ref{lem:mobius-cancel} to
\(n=\sum_{d\mid n}\varphi(d)\), obtaining the first formula.  Substitute
\(d\leftrightarrow n/d\) for the second.
\end{proof}

\begin{theorem}[Summatory floor-square identity]\label{thm:summatory-phi}
For every \(N\ge1\),
\[
 \boxed{\displaystyle
 2\sum_{n=1}^{N}\varphi(n)-1
 =\sum_{d=1}^{N}\mu(d)\left\lfloor\frac Nd\right\rfloor^2.}
\]
\end{theorem}

\begin{proof}
Let \(C_N\) be the number of ordered pairs \((a,b)\) with
\(1\le a,b\le N\) and \(\gcd(a,b)=1\).  By M\"obius cancellation,
\[
 \ind{\gcd(a,b)=1}=\sum_{d\mid a,\ d\mid b}\mu(d).
\]
Summing over the finite square gives
\[
 C_N=\sum_{d\le N}\mu(d)\#\{a\le N:d\mid a\}^2
 =\sum_{d\le N}\mu(d)\lfloor N/d\rfloor^2.
\]
On the other hand, the diagonal contributes only \((1,1)\).  Above the
diagonal, for each \(b=2,\ldots,N\), there are \(\varphi(b)\) choices of
\(a<b\) coprime to \(b\); reflection supplies the same number below the
diagonal.  Thus
\(C_N=1+2\sum_{b=2}^N\varphi(b)=2\sum_{b=1}^N\varphi(b)-1\).
All exchanges are bounded sums.
\end{proof}

\begin{corollary}[Farey sequence count]
The number of reduced fractions in \([0,1]\) with denominator at most \(N\)
is
\[
 1+\sum_{n=1}^{N}\varphi(n).
\]
\end{corollary}

\begin{proof}
Include \(0/1\).  For every denominator \(n\), the reduced fractions
\(a/n\) with \(1\le a\le n\) are counted by \(\varphi(n)\), including
\(1/1\) at \(n=1\).  Uniqueness follows from cross multiplication and
coprimality.  This is a bounded enumeration theorem.
\end{proof}

\section{A calculus of totient identities}

The following identities are most transparently proved by comparing the local
prime-power factors from Theorem~\ref{thm:phi-product}.  Such comparison is
constructive because the factor vector is explicitly computed.

\begin{theorem}[Power identity]\label{thm:phi-power}
For \(n\ge1\) and \(k\ge1\),
\[
 \boxed{\varphi(n^k)=n^{k-1}\varphi(n).}
\]
\end{theorem}

\begin{proof}
If \(n=\prod p^{a_p}\), then \(n^k=\prod p^{ka_p}\), so
\[
 \varphi(n^k)=\prod p^{ka_p-1}(p-1)
 =\left(\prod p^{(k-1)a_p}\right)
   \left(\prod p^{a_p-1}(p-1)\right).
\]
The two factors are \(n^{k-1}\) and \(\varphi(n)\).
\end{proof}

\begin{theorem}[Divisibility monotonicity]\label{thm:phi-divides}
If \(a\mid b\) and \(a,b>0\), then \(\varphi(a)\mid\varphi(b)\).
\end{theorem}

\begin{proof}
Write the prime exponents as \(\alpha_p\le\beta_p\).  For a prime already
present in \(a\), the ratio of local factors is \(p^{\beta_p-\alpha_p}\).
For a new prime of \(b\), its entire local factor
\(p^{\beta_p-1}(p-1)\) is an integer factor.  Multiplying these finitely many
ratios gives an explicit quotient \(Q\) satisfying
\(\varphi(b)=Q\varphi(a)\).
\end{proof}

\begin{theorem}[Noncoprime multiplication formula]\label{thm:phi-general-mul}
Let \(d=\gcd(m,n)>0\).  Then
\[
 \boxed{\varphi(mn)\varphi(d)=\varphi(m)\varphi(n)d.}
\]
Equivalently, in ordinary rational notation,
\(
 \varphi(mn)=\varphi(m)\varphi(n)d/\varphi(d).
\)
\end{theorem}

\begin{proof}
Fix a prime and write \(a=v_p(m), b=v_p(n)\), with
\(r=\min(a,b)\).  If one exponent is zero, the local equality is immediate.
If both are positive, the exponents of \(p\) and the two factors \((p-1)\)
are
\[
 [a+b-1]+[r-1]
 \quad\text{on the left, and}\quad
 [a-1]+[b-1]+r
 \quad\text{on the right},
\]
which are equal.  Multiplying the local equalities proves the result.
\end{proof}

\begin{corollary}[Gcd-lcm identity]\label{cor:phi-gcd-lcm}
For positive \(m,n\),
\[
 \boxed{\varphi(\gcd(m,n))\varphi(\operatorname{lcm}(m,n))
 =\varphi(m)\varphi(n).}
\]
\end{corollary}

\begin{proof}
Primewise, the gcd and lcm exponents are the minimum and maximum of the two
exponents.  Their multiset is the same as the original pair, so the product of
local totient factors is unchanged.
\end{proof}

\begin{theorem}[Doubling]\label{thm:phi-double}
For \(n\ge1\),
\[
 \varphi(2n)=\begin{cases}
 \varphi(n),&n\text{ odd},\\
 2\varphi(n),&n\text{ even}.
 \end{cases}
\]
\end{theorem}

\begin{proof}
If \(n\) is odd, multiplicativity with \(\varphi(2)=1\) applies.  If
\(n=2^a m\) with \(a\ge1\) and \(m\) odd, the local factor at two changes
from \(2^{a-1}\) to \(2^a\), while every odd local factor is unchanged.
\end{proof}

\begin{theorem}[Radical identity]\label{thm:phi-rad}
For \(n>0\),
\[
 \boxed{\varphi(n)\rad(n)=n\varphi(\rad(n)).}
\]
\end{theorem}

\begin{proof}
For each \(p^a\Vert n\), the local factor on the left is
\(p^{a-1}(p-1)\cdot p=p^a(p-1)\), equal to the local factor on the right.
\end{proof}

\begin{theorem}[Squarefree-divisor identity]\label{thm:squarefree-divisor}
For \(n>0\),
\[
 \boxed{\displaystyle
 \sum_{d\mid n}\frac{\mu(d)^2}{\varphi(d)}=\frac n{\varphi(n)}.}
\]
Internally this denotes the equality obtained after multiplying by the product
of the finitely many nonzero denominators.
\end{theorem}

\begin{proof}
Only squarefree divisors contribute.  Finite distributivity gives
\[
 \sum_{d\mid n}\frac{\mu(d)^2}{\varphi(d)}
 =\prod_{p\mid n}\left(1+\frac1{p-1}\right)
 =\prod_{p\mid n}\frac p{p-1}
 =\frac n{\varphi(n)}.
\]
Every step can be cross-multiplied prime by prime, so no rational object is
needed internally.
\end{proof}

\section{Parity, reduced residues, and Menon's identity}

\begin{theorem}[Evenness and powers of two]\label{thm:phi-even}
If \(n>2\), then \(\varphi(n)\) is even.  More generally, if \(n\) has
\(r\) distinct odd prime divisors, then \(2^r\mid\varphi(n)\).
\end{theorem}

\begin{proof}
For evenness, pair each reduced residue \(a\) with \(n-a\).  Both are reduced,
and they are equal modulo \(n\) only if \(2a\equiv0\pmod n\).  Since \(a\)
is a unit, this would give \(n\mid2\), impossible for \(n>2\).  Hence the
reduced residue set is partitioned into two-element orbits.

For the stronger statement, each distinct odd prime divisor \(p\) contributes
the even factor \(p-1\) to the product formula.  The remaining factors are
integers, so their product contains at least \(r\) factors of two.
\end{proof}

\begin{theorem}[Sum of the totatives]\label{thm:sum-totatives}
For \(n>1\), if
\(S(n)=\sum_{1\le a<n,\ \gcd(a,n)=1}a\), then
\[
 \boxed{2S(n)=n\varphi(n).}
\]
Thus \(S(n)=n\varphi(n)/2\) in ordinary notation.
\end{theorem}

\begin{proof}
For \(n>2\), use the fixed-point-free involution
\(a\mapsto n-a\); every pair sums to \(n\), and there are
\(\varphi(n)/2\) pairs.  For \(n=2\), the sole totative is one and the
cross-multiplied identity reads \(2=2\).
\end{proof}

\begin{theorem}[Menon's identity]\label{thm:menon}
For every positive \(n\),
\[
 \boxed{\displaystyle
 \sum_{\substack{1\le a\le n\\\gcd(a,n)=1}}\gcd(a-1,n)
 =\varphi(n)\tau(n),}
\]
where \(\tau(n)\) is the number of positive divisors of \(n\).
\end{theorem}

\begin{proof}
Use Theorem~\ref{thm:sum-phi} in the form
\(\gcd(a-1,n)=\sum_{d\mid\gcd(a-1,n)}\varphi(d)\).  Swapping bounded sums gives
\[
 \sum_{d\mid n}\varphi(d)N_d,
 \qquad
 N_d=\#\{a\bmod n:\gcd(a,n)=1,\ a\equiv1\pmod d\}.
\]
We show \(N_d=\varphi(n)/\varphi(d)\) by an explicit local count.  Write
\(n=\prod p^{\alpha_p}\) and
\(d=\prod p^{\delta_p}\), with
\(0\le\delta_p\le\alpha_p\).  CRT makes the count a product over primes.
For one prime power:
\begin{itemize}
\item if \(\delta=0\), all \(\varphi(p^\alpha)\) units are allowed;
\item if \(\delta>0\), the residues congruent to one modulo \(p^\delta\)
are \(1+p^\delta t\) with \(0\le t<p^{\alpha-\delta}\); all are units, so
there are \(p^{\alpha-\delta}\) choices.
\end{itemize}
The quotient of these local counts by \(\varphi(p^\alpha)\) is exactly
\(1/\varphi(p^\delta)\), interpreted by the explicit integer quotient.  Thus
\(\varphi(d)N_d=\varphi(n)\) for every divisor \(d\).  There are \(\tau(n)\)
divisors, and the result follows.
\end{proof}

\section{Euler's theorem, orders, and RSA}

\begin{theorem}[Euler's theorem]\label{thm:euler}
If \(n>0\) and \(\gcd(a,n)=1\), then
\[
 \boxed{a^{\varphi(n)}\equiv1\pmod n.}
\]
For prime \(p\), this gives Fermat's little theorem
\(a^{p-1}\equiv1\pmod p\) whenever \(p\nmid a\).
\end{theorem}

\begin{proof}
Let \(R\) be the bounded list of reduced residues modulo \(n\).  Multiplication
by \(a\), followed by canonical remainder, is a permutation of \(R\): its
inverse is multiplication by the Bezout inverse of \(a\).  Therefore
\[
 \prod_{r\in R}r\equiv\prod_{r\in R}ar
 =a^{\varphi(n)}\prod_{r\in R}r\pmod n.
\]
For \(n=1\) the congruence is immediate.  For \(n>1\), the product of the elements of \(R\) is a unit modulo \(n\), so the extended
Euclidean algorithm computes an inverse and permits cancellation.  This gives
the desired congruence.  The permutation and product identities are finite
bounded inductions.
\end{proof}

\begin{corollary}[Order divides the totient]\label{cor:order-divides}
For \(\gcd(a,n)=1\), let \(r\) be the least positive exponent with
\(a^r\equiv1\pmod n\).  Then \(r\mid\varphi(n)\).
\end{corollary}

\begin{proof}
Euler's theorem supplies such an exponent.  Divide
\(\varphi(n)=qr+s\) with \(s<r\).  Then
\(1\equiv a^{\varphi(n)}\equiv(a^r)^q a^s\equiv a^s\).  Minimality of
\(r\) forces \(s=0\).  Least positive search is bounded by \(\varphi(n)\),
so the order function is primitive recursive on units.
\end{proof}

\begin{theorem}[Zsigmondy-type divisibility]\label{thm:order-am1}
If \(a>1\) and \(m\ge1\), then
\[
 \boxed{m\mid\varphi(a^m-1).}
\]
\end{theorem}

\begin{proof}
Let \(N=a^m-1\).  Since \(a\) and \(N\) are coprime,
\(a\) is a unit modulo \(N\), and \(a^m\equiv1\pmod N\).  If
\(0<r<m\) also satisfied \(a^r\equiv1\pmod N\), then
\(N\mid a^r-1\), but
\(0<a^r-1<a^m-1=N\), impossible.  Hence the order of \(a\) modulo \(N\) is
exactly \(m\), and Corollary~\ref{cor:order-divides} gives the result.
\end{proof}

\begin{theorem}[RSA correctness]\label{thm:rsa}
Let \(p\ne q\) be primes, \(N=pq\), and suppose
\(ed\equiv1\pmod{\varphi(N)}\).  Then for every message \(x\),
\[
 (x^e)^d\equiv x\pmod N.
\]
\end{theorem}

\begin{proof}
Write \(ed=1+t\varphi(N)\).  Modulo \(p\), either \(p\mid x\), in which case
both sides are zero, or \(x\) is a unit.  Since
\(p-1\mid\varphi(pq)=(p-1)(q-1)\), Fermat's theorem gives
\(x^{t\varphi(N)}\equiv1\pmod p\), hence \(x^{ed}\equiv x\pmod p\).
The same holds modulo \(q\).  The constructive CRT then combines the two
congruences into one modulo \(pq\).  The proof handles nonunits explicitly and
uses no classical group theory.
\end{proof}


\section{Further finite statements from the totient reference pages}

\begin{theorem}[Subbarao congruence classification]\label{thm:subbarao}
For every prime \(p\),
\[
 p\sigma(p)\equiv2\pmod{\varphi(p)}.
\]
The only composite positive integers satisfying
\[
 n\sigma(n)\equiv2\pmod{\varphi(n)}
\]
are \(4,6,22\).
\end{theorem}

\begin{proof}
For a prime \(p\), \(\sigma(p)=p+1\) and \(\varphi(p)=p-1\).  Since
\(p\equiv1\pmod{p-1}\), one has
\(p(p+1)\equiv1\cdot2=2\).

Now let composite \(n\) satisfy the congruence.  If an odd prime square
\(q^2\) divides \(n\), then \(q\mid\varphi(n)\) and \(q\mid n\), so the
congruence modulo \(q\) reads \(0\equiv2\), impossible.  If \(8\mid n\),
then \(4\mid\varphi(n)\) and \(4\mid n\sigma(n)\), again impossible.  If
\(4\Vert n\), write \(n=4m\) with \(m\) odd.  When \(m>1\), evenness of
\(\varphi(m)\) gives \(4\mid\varphi(4m)=2\varphi(m)\), while
\(4\mid n\sigma(n)\), a contradiction.  Therefore this case yields only
\(n=4\), which directly satisfies the congruence.

All remaining composite candidates are squarefree.  If two distinct odd
primes \(p,q\) divide \(n\), then \((p-1)(q-1)\mid\varphi(n)\), so
\(4\mid\varphi(n)\); also \((p+1)(q+1)\mid\sigma(n)\), so
\(4\mid n\sigma(n)\).  This again contradicts a residue of two modulo four.
Thus a composite candidate must be \(n=2p\) for one odd prime \(p\).  Here
\[
 \varphi(2p)=p-1,
 \qquad
 \sigma(2p)=3(p+1),
\]
so modulo \(p-1\),
\[
 n\sigma(n)=6p(p+1)\equiv12.
\]
The required congruence is equivalent to \(p-1\mid10\).  Since \(p\) is an
odd prime, this gives \(p=3\) or \(p=11\), hence \(n=6\) or \(22\).  Direct
calculation verifies both.  Every step is a factorization case distinction
on bounded data and therefore formalizable in \(\ipra\).
\end{proof}

\begin{theorem}[The \(a^n+b^n\) totient congruence]\label{thm:plus-power}
For distinct positive integers \(a,b\) and every \(n\ge1\),
\[
 \boxed{n\mid\varphi(a^n+b^n).}
\]
In fact, after dividing \(a,b\) by their gcd, the unit constructed below has
order \(n\) when \(n\) is odd and order \(2n\) when \(n\) is even.
\end{theorem}

\begin{proof}
Let \(g=\gcd(a,b)\), \(A=a/g\), \(B=b/g\), and
\(C=A^n+B^n\).  Then \(\gcd(A,B)=1\), \(A\ne B\), and
\(C\mid a^n+b^n=g^nC\).  By Theorem~\ref{thm:phi-divides}, it is enough to
show \(n\mid\varphi(C)\).  Since
\(\gcd(B,C)=\gcd(B,A^n)=1\), a Bezout inverse \(B^{-1}\pmod C\) is
primitive-recursively available.

Suppose first that \(n\) is odd and put
\(u\equiv-A B^{-1}\pmod C\).  The relation
\(A^n\equiv-B^n\pmod C\) gives \(u^n\equiv1\).  If
\(0<k<n\) and \(u^k\equiv1\), then, after multiplying by \(B^k\), either
\(C\mid A^k-B^k\) or \(C\mid A^k+B^k\), according to the parity of \(k\).
The relevant nonzero absolute value is strictly smaller than
\(A^n+B^n=C\), impossible.  Hence \(u\) has order exactly \(n\).

If \(n\) is even, put \(u\equiv A B^{-1}\pmod C\).  Then
\(u^n\equiv-1\) and \(u^{2n}\equiv1\).  No exponent
\(0<k<n\) can give one, because that would make \(C\) divide the nonzero
number \(A^k-B^k\) of smaller absolute value.  At \(k=n\) the value is
minus one; note that \(C\ge5\).  If \(n<k<2n\), write \(k=n+r\) with
\(0<r<n\); then \(u^k=1\) would imply \(u^r=-1\), so
\(C\mid A^r+B^r\), again impossible by size.  Thus the order is exactly
\(2n\).

In either case, the order divides \(\varphi(C)\) by
Corollary~\ref{cor:order-divides}; hence \(n\mid\varphi(C)\), and
\(\varphi(C)\mid\varphi(a^n+b^n)\) completes the proof.  The argument is
entirely finite and avoids using Zsigmondy's theorem as a black box.
\end{proof}

\begin{proposition}[Goldbach implication]\label{prop:goldbach-phi}
Assume the strong Goldbach statement that every even integer at least four is
a sum of two primes.  Then for every positive \(m\) there are primes \(p,q\)
with
\[
 \varphi(p)+\varphi(q)=2m.
\]
\end{proposition}

\begin{proof}
Apply Goldbach to \(2m+2\), obtaining primes \(p+q=2m+2\).  Then
\(\varphi(p)+\varphi(q)=(p-1)+(q-1)=2m\).  This is a constructive implication:
a witness-producing Goldbach procedure immediately produces the totient
witnesses.  It does not prove Goldbach itself.
\end{proof}

The congruence \(n-1\equiv0\pmod{\varphi(n)}\) is automatic for primes, but no
composite solution is known; this is Lehmer's totient problem and remains
status \StatusO.  Equal-totient examples, the unique reported triple below
\(10^{10}\), arithmetic progressions of equal totients, and examples of
\(\varphi(\sigma(n))=n\) are finite calculations.  They are reproduced by the
archive's exact-integer checker rather than promoted to universal theorems.

\section{Cyclotomic and generating-function identities}

\begin{theorem}[Cyclotomic polynomial congruence]\label{thm:cyclotomic-gcd-sum}
Define
\[
 P_n(X)=\sum_{k=1}^{n}\gcd(k,n)X^k\in\Z[X].
\]
For \(n\ge1\),
\[
 \boxed{P_n(X)\equiv\varphi(n)\pmod{\Phi_n(X)}.}
\]
Consequently, for every primitive \(n\)-th root of unity \(\zeta\),
\[
 \sum_{k=1}^{n}\gcd(k,n)\zeta^k=\varphi(n),
\]
and taking real parts yields the cosine identity on the reference page.
\end{theorem}

\begin{proof}
In the quotient ring \(\Z[X]/(\Phi_n)\), let \(z\) be the class of \(X\).
It satisfies \(z^n=1\), and for every proper divisor \(d<n\) of \(n\),
\(z^d\ne1\) in the universal cyclotomic extension.  Use the divisor-sum
identity on each coefficient:
\[
 \gcd(k,n)=\sum_{d\mid\gcd(k,n)}\varphi(d).
\]
Then
\[
 P_n(z)=\sum_{d\mid n}\varphi(d)
          \sum_{1\le j\le n/d}z^{dj}.
\]
For \(d<n\), write \(m=n/d>1\).  The inner polynomial is
\[
 X^d+X^{2d}+\cdots+X^{md}
 =X^d\frac{X^n-1}{X^d-1}.
\]
Because \(d\mid n\) is proper, the cyclotomic factor \(\Phi_n(X)\) occurs in
\((X^n-1)/(X^d-1)\), so this entire inner polynomial is zero modulo
\(\Phi_n\).  For \(d=n\), the inner polynomial is \(X^n\), congruent to one
modulo \(\Phi_n\).  Thus only \(\varphi(n)\) remains.

For a literal \(\ipra\) formulation, compute \(\Phi_n\) by the finite recursion
\(X^n-1=\prod_{d\mid n}\Phi_d(X)\), perform signed polynomial remainder by
bounded coefficient arrays, and assert that the remainder of
\(P_n-\varphi(n)\) is zero.  Each degree is at most \(n\), so all loops are
bounded and primitive recursive.  Evaluation at a complex root is the
external algebraic corollary.
\end{proof}

\begin{theorem}[Lambert-series coefficient identity]\label{thm:lambert}
As a formal power-series identity,
\[
 \boxed{\displaystyle
 \sum_{n\ge1}\frac{\varphi(n)q^n}{1-q^n}=\frac q{(1-q)^2}.}
\]
Precisely, for every \(m\ge1\), the coefficient of \(q^m\) on each side is
\(m\), and this coefficient equality is a theorem of \(\ipra\).
\end{theorem}

\begin{proof}
The coefficient of \(q^m\) in
\(q^n/(1-q^n)=q^n+q^{2n}+\cdots\) is one exactly when \(n\mid m\).  Therefore
the left coefficient is \(\sum_{n\mid m}\varphi(n)=m\).  The formal identity
\(q/(1-q)^2=\sum_{m\ge1}mq^m\) follows by a finite induction on coefficients,
or from the product of the two finite geometric truncations with a remainder
above degree \(m\).  Each fixed coefficient depends on only finitely many
terms.
\end{proof}

\begin{theorem}[Formal Dirichlet-series identity]\label{thm:phi-dirichlet}
Let \(\mathbf{id}(n)=n\), let \(\mathbf{1}(n)=1\), and use finite Dirichlet
convolution
\((f*g)(n)=\sum_{d\mid n}f(d)g(n/d)\).  Then
\[
 \boxed{\varphi=\mathbf{id}*\mu,\qquad
        \varphi*\mathbf{1}=\mathbf{id}.}
\]
The analytic notation
\[
 \sum_{n\ge1}\frac{\varphi(n)}{n^s}
 =\frac{\zeta(s-1)}{\zeta(s)}
\]
is the convergent-series interpretation of this exact coefficient theorem.
\end{theorem}

\begin{proof}
The first equality is Corollary~\ref{cor:phi-mobius}; the second is
Theorem~\ref{thm:sum-phi}.  Associativity and commutativity of convolution are
proved by reindexing a finite divisor triangle.  Since
\(\mu*\mathbf1=\varepsilon\), the identities are equivalent.  No convergence
claim is part of the \(\ipra\) theorem.
\end{proof}

\section{Elementary bounds and extremal witness forms}

\begin{theorem}[Elementary lower bound]\label{thm:phi-lower}
For every positive \(n\),
\[
 \boxed{2\varphi(n)^2\ge n.}
\]
Moreover,
\[
 \varphi(n)^2\ge n
 \quad\text{for every }n\notin\{2,6\}.
\]
The square-root statements on the reference pages follow externally because
all quantities are nonnegative.
\end{theorem}

\begin{proof}
From the prime-power formula,
\[
 \frac{\varphi(p^a)^2}{p^a}=p^{a-2}(p-1)^2.
\]
For \(p=2\), this is \(2^{a-2}\), at least \(1/2\).  For every odd prime,
it is at least \((p-1)^2/p>1\).  Multiplying local inequalities proves
\(\varphi(n)^2/n\ge1/2\).

For the sharper assertion, a local ratio below one can occur only at the
factor \(2^1\), whose ratio is \(1/2\).  If an odd prime \(p\ge5\) occurs, its
local ratio is at least \(16/5>2\), overcoming that loss.  If a power
\(3^a\) occurs, its local ratio is \(4/3\) for \(a=1\) and at least \(4\)
for \(a\ge2\).  Thus the only products below one are \(n=2\) and
\(n=2\cdot3=6\); the boundary case \(n=4\) has ratio one.  These comparisons
are integer cross multiplications.
\end{proof}

\begin{theorem}[Composite upper bound]\label{thm:phi-upper-composite}
If \(n>1\) is composite, then
\[
 \boxed{\varphi(n)\le n-\lceil\sqrt n\rceil}.
\]
Equivalently, if \(r\) is the least positive integer satisfying \(n\le r^2\),
then \(\varphi(n)\le n-r\).  In particular, for \(n>1\),
\(\varphi(n)\le n-1\), with equality exactly when \(n\) is prime.
\end{theorem}

\begin{proof}
Let \(p\) be the least prime divisor of composite \(n\).  Then
\(p\le\sqrt n\), for otherwise the complementary factor \(n/p\) would be a
smaller divisor greater than one.  At least the \(n/p\) multiples of \(p\)
below \(n\) are excluded from the totatives, so
\(\varphi(n)\le n-n/p\le n-\sqrt n\).  Replace square roots by the displayed
integer bound.  For a prime, Theorem~\ref{thm:phi-prime-power} gives equality
\(n-1\); for a composite, the stronger bound is strict below \(n-1\).
\end{proof}

\begin{theorem}[Totient-sigma product]\label{thm:phi-sigma}
Let \(\sigma(n)=\sum_{d\mid n}d\).  For \(n>1\),
\[
 \boxed{\displaystyle
 \frac{\varphi(n)\sigma(n)}{n^2}
 =\prod_{p^a\Vert n}\left(1-\frac1{p^{a+1}}\right),
 \qquad
 \frac12<\frac{\varphi(n)\sigma(n)}{n^2}<1.}
\]
Both inequalities have denominator-cleared \(\ipra\) forms.
\end{theorem}

\begin{proof}
For a prime power,
\(
 \sigma(p^a)=(p^{a+1}-1)/(p-1)
\), so multiplication by
\(\varphi(p^a)=p^{a-1}(p-1)\) and division by \(p^{2a}\) gives the local
factor.  Multiplicativity gives the product.  Every local factor is below one.
For the lower bound,
\[
 \prod_{p^a\Vert n}\left(1-p^{-a-1}\right)
 \ge\prod_{p\mid n}(1-p^{-2})
 \ge\prod_{m=2}^{n}(1-m^{-2}).
\]
The last finite product telescopes:
\[
 \prod_{m=2}^{n}\frac{(m-1)(m+1)}{m^2}
 =\frac{n+1}{2n}>\frac12.
\]
The telescoping product is already strictly greater than one half, so no
strictness is needed in the preceding subset comparison.  All displayed
rational inequalities are cross-multiplied finite products.
\end{proof}

The sharper classical constant \(6/\pi^2\) follows by identifying the infinite
Euler product \(\prod_p(1-p^{-2})\).  The present theorem records the exact
finite identity and a uniform rational constant provable in native \(\ipra\);
the \(6/\pi^2\) refinement is status \StatusA because it requires a separate
constructive development of \(\pi\), the Basel sum, and an infinite-product
limit.

\begin{theorem}[Finite limsup witnesses]\label{thm:limsup-phi}
For every \(K\ge2\), one can primitive-recursively find \(n>K\) such that
\[
 K\varphi(n)\ge(K-1)n.
\]
Thus the usual statement \(\limsup_{n\to\infty}\varphi(n)/n=1\) has an exact
finite witness form in \(\ipra\).
\end{theorem}

\begin{proof}
Euclid's proof is an explicit bounded construction: take a prime divisor
\(p\) of \(K!+1\).  Then \(p>K\).  Set \(n=p\).  Since
\(\varphi(p)=p-1\),
\[
 K\varphi(p)=K(p-1)=(K-1)p+(p-K)\ge(K-1)p.
\]
A least prime divisor is obtained by bounded search.
\end{proof}

\begin{theorem}[Finite liminf witnesses]\label{thm:liminf-phi}
For every \(K\ge1\), there is a primitive-recursively specified positive
integer \(n\) such that
\[
 K\varphi(n)<n.
\]
Hence the finite content of
\(\liminf_{n\to\infty}\varphi(n)/n=0\) is provable in \(\ipra\).
\end{theorem}

\begin{proof}
Let \(n\) be the product of all primes at most a bound \(B\).  Then
\[
 \frac n{\varphi(n)}=\prod_{p\le B}\frac p{p-1}.
\]
The finite Euler-product argument used in the constructive proof of divergence
of reciprocal primes shows, with an explicit primitive-recursive witness
\(B=B(K)\), that this product exceeds \(K\).  One elementary choice follows
from harmonic domination:
\[
 H_B\le\prod_{p\le B}\frac p{p-1},
\]
and a power-of-two block proof gives \(H_{2^{2K}}>K\).  Both inequalities are
finite, denominator-cleared statements.  Taking the primorial up to
\(B=2^{2K}\) therefore gives \(n/\varphi(n)>K\), which is the desired
cross-multiplied inequality.  The factorization injection and finite Euler
product are developed in the companion ProveIt prime-reciprocal package and
can be reused verbatim.
\end{proof}

\section{Nontotients and finite classification results}

\begin{theorem}[Odd nontotients]\label{thm:odd-nontotient}
No odd integer greater than one is a value of \(\varphi\).
\end{theorem}

\begin{proof}
The values at one and two are one.  For every input greater than two,
Theorem~\ref{thm:phi-even} says the value is even.  Hence an odd value greater
than one never occurs.
\end{proof}

\begin{theorem}[Solutions of \(\varphi(n)=2\)]\label{thm:phi-two}
\[
 \boxed{\varphi(n)=2\quad\Longleftrightarrow\quad n\in\{3,4,6\}.}
\]
\end{theorem}

\begin{proof}
Suppose \(\varphi(n)=2\).  Any odd prime divisor \(p\) contributes the factor
\(p-1\), so \(p-1\mid2\), forcing \(p=3\).  Its exponent can only be one,
because \(\varphi(3^a)=2\cdot3^{a-1}\).  The power of two dividing \(n\)
can have exponent at most two; its local factor is \(2^{b-1}\) for \(b\ge1\).
If no odd prime occurs, equality forces \(b=2\), giving \(n=4\).  If the
factor three occurs, the two-local factor must be one, so \(b=0\) or \(1\),
giving \(3\) or \(6\).  Direct calculation verifies all three.
\end{proof}

\begin{proposition}[Arithmetic part of the constructible prime-power criterion]
Let \(p\) be an odd prime and \(a\ge1\).  If \(\varphi(p^a)\) is a power of
two, then \(a=1\) and \(p-1\) is a power of two.
\end{proposition}

\begin{proof}
The factorization
\(\varphi(p^a)=p^{a-1}(p-1)\) contains an odd factor \(p^{a-1}\).  A power
of two has no odd factor greater than one, so \(a-1=0\).  The remaining factor
is \(p-1\), which must itself be a power of two.  The further conclusion that
such a prime is a Fermat prime uses the elementary factorization of
\(2^m+1\) when \(m\) has an odd divisor and is likewise formalizable in
\(\ipra\).
\end{proof}

Carmichael's conjecture that no totient value occurs exactly once is not proved
here and remains status \StatusO.  The finite checker in the archive verifies
many equal-totient examples from the reference pages but, correctly, makes no
claim about the conjecture.

\section{The weighted von Mangoldt function}

\begin{definition}[Prime-power recognizer]\label{def:primepow}
For \(n\ge2\), define \(\operatorname{ppBase}(n)\) to be the least prime
\(p\le n\) for which there is an exponent \(1\le a\le n\) with \(p^a=n\),
and return zero if no such pair exists.  Define
\(\operatorname{ppExp}(n)\) as the least such exponent for the selected base.
Both are primitive recursive bounded searches.
\end{definition}

Uniqueness follows from valuations: if \(p^a=q^b>1\) with primes \(p,q\),
then applying \(v_p\) shows \(p=q\), and then strict monotonicity of powers
shows \(a=b\).

\begin{definition}[Weighted von Mangoldt function]\label{def:weighted-lambda}
For a primitive-recursive prime weight \(w\), put
\[
 \Lambda_w(n)=\begin{cases}
 w(p),&n=p^a\text{ for a prime }p\text{ and }a\ge1,\\
 0,&\text{otherwise}.
 \end{cases}
\]
Thus \(\Lambda_w(1)=0\).  The classical function is the specialization
\(w(p)=\log p\).
\end{definition}

For example, the support begins at
\(2,3,4,5,7,8,9,11,13,16,17,19,23,25,27,\ldots\), and the value at every
power of the same prime is the same weight.  The function is neither
multiplicative nor additive: for distinct primes \(p,q\) with nonzero weights,
\(\Lambda_w(pq)=0\), while both \(\Lambda_w(p)\) and
\(\Lambda_w(q)\) are nonzero.

\section{Mangoldt divisor sums and M\"obius identities}

\begin{theorem}[Fundamental divisor sum]\label{thm:lambda-divsum}
For every positive \(n\),
\[
 \boxed{\displaystyle
 L_w(n)=\sum_{d\mid n}\Lambda_w(d).}
\]
For \(w(p)=\log p\), this is
\(\sum_{d\mid n}\Lambda(d)=\log n\).
\end{theorem}

\begin{proof}
Write \(n=\prod_pp^{a_p}\).  The divisors on which \(\Lambda_w\) is nonzero
are exactly \(p,p^2,\ldots,p^{a_p}\) for the primes dividing \(n\).  Each of
those \(a_p\) terms contributes \(w(p)\).  Therefore the divisor sum is
\(\sum_p a_pw(p)=L_w(n)\).  This is a finite reindexing of the computed factor
vector.
\end{proof}

\begin{theorem}[M\"obius forms of Mangoldt]\label{thm:lambda-mobius}
For \(n>0\),
\[
 \boxed{\displaystyle
 \Lambda_w(n)=\sum_{d\mid n}\mu(d)L_w(n/d)
 =\sum_{d\mid n}\mu(n/d)L_w(d).}
\]
For \(n>1\), also
\[
 \boxed{\displaystyle
 \Lambda_w(n)=-\sum_{d\mid n}\mu(d)L_w(d).}
\]
\end{theorem}

\begin{proof}
The first equality is M\"obius inversion applied to
Theorem~\ref{thm:lambda-divsum}.  The second is the divisor substitution
\(d\leftrightarrow n/d\).  For the third, complete additivity gives
\(L_w(n/d)=L_w(n)-L_w(d)\) in signed arithmetic.  Hence
\[
 \Lambda_w(n)
 =L_w(n)\sum_{d\mid n}\mu(d)-\sum_{d\mid n}\mu(d)L_w(d).
\]
For \(n>1\), the first sum is zero by Lemma~\ref{lem:mobius-cancel}.
\end{proof}

\begin{theorem}[M\"obius--Mangoldt convolution]\label{thm:mu-lambda}
For positive \(n\),
\[
 \boxed{\displaystyle
 \sum_{d\mid n}\mu(n/d)\Lambda_w(d)=-\mu(n)L_w(n).}
\]
\end{theorem}

\begin{proof}
It suffices to inspect the prime factorization.
If \(n\) contains at least two distinct primes and is not squarefree, both
sides are zero after cancellation; the finite proof below treats all cases
uniformly.  Substitute the first formula of
Theorem~\ref{thm:lambda-mobius}:
\[
 \sum_{d\mid n}\mu(n/d)\sum_{e\mid d}\mu(e)L_w(d/e).
\]
Reindex by \(d=er\) and use finite convolution associativity.  In convolution
notation the left side is \(\mu*\Lambda_w=\mu*(\mu*L_w)\).  A direct local
calculation on \(p^a\) shows that the convolution is zero unless \(n\) is
squarefree; for squarefree \(n=p_1\cdots p_r\), only divisors \(d=p_i\)
contribute to \(\Lambda_w(d)\), and
\[
 \sum_i\mu(n/p_i)w(p_i)=(-1)^{r-1}\sum_iw(p_i)
 =-\mu(n)L_w(n).
\]
The nonsquarefree case is similarly paired by toggling one repeated-prime
factor.  This pairing is a bounded bit-vector involution and formalizes by
quantifier-free induction on the prime list.
\end{proof}

The three theorems above contain the standard formulas displayed by both
Mangoldt reference pages.  They are exact \(\ipra\) theorems for every
primitive-recursive weight, not merely identities of real-valued functions.

\section{Least common multiples and the Chebyshev function}

\begin{definition}
Let
\[
 \lcmu(N)=\operatorname{lcm}(1,2,\ldots,N),
 \qquad \lcmu(0)=1.
\]
This is a primitive-recursive bounded fold using
\(\operatorname{lcm}(a,b)=ab/\gcd(a,b)\), with exact division.
\end{definition}

\begin{lemma}[Valuation of the initial lcm]\label{lem:lcm-valuation}
For prime \(p\),
\[
 v_p(\lcmu(N))=\max\{a:p^a\le N\}.
\]
\end{lemma}

\begin{proof}
The valuation of an lcm is the maximum of the valuations.  Among
\(1,\ldots,N\), an exponent at least \(a\) occurs exactly when \(p^a\le N\).
Both maxima are bounded by \(N\), and the equality follows by induction over
the lcm fold.
\end{proof}

\begin{theorem}[Lcm quotient]\label{thm:lcm-quotient}
For \(N\ge1\),
\[
 \frac{\lcmu(N)}{\lcmu(N-1)}=
 \begin{cases}
 p,&N=p^a\text{ for a prime }p,\ a\ge1,\\
 1,&\text{otherwise}.
 \end{cases}
\]
The quotient is exact.  Consequently
\[
 L_w(\lcmu(N))-L_w(\lcmu(N-1))=\Lambda_w(N).
\]
For logarithmic weight this is the MathWorld formula
\(e^{\Lambda(N)}=\lcmu(N)/\lcmu(N-1)\).
\end{theorem}

\begin{proof}
By Lemma~\ref{lem:lcm-valuation}, moving from \(N-1\) to \(N\) increases a
prime valuation exactly when \(N\) itself is a new power \(p^a\); then it
increases that single valuation by one.  Two different prime-power
representations are impossible.  Thus the quotient is either one or the
corresponding prime.  Complete additivity of \(L_w\) gives the difference
formula.
\end{proof}

\begin{definition}[Weighted Chebyshev sum]
\[
 \Psi_w(N)=\sum_{1\le n\le N}\Lambda_w(n).
\]
\end{definition}

\begin{theorem}[Chebyshev-lcm identity]\label{thm:psi-lcm}
For every \(N\),
\[
 \boxed{\displaystyle
 \Psi_w(N)=\sum_{p^a\le N}w(p)=L_w(\lcmu(N)).}
\]
For \(w(p)=\log p\), this is
\(\psi(N)=\log(\operatorname{lcm}(1,\ldots,N))\).
\end{theorem}

\begin{proof}
The first equality merely expands the support of \(\Lambda_w\).  Summing the
successive-difference identity of Theorem~\ref{thm:lcm-quotient} telescopes:
\[
 \sum_{n=1}^{N}\Lambda_w(n)
 =L_w(\lcmu(N))-L_w(\lcmu(0))=L_w(\lcmu(N)),
\]
because \(\lcmu(0)=1\) and \(L_w(1)=0\).  Telescoping is a primitive
recursion identity.
\end{proof}

\section{Formal Dirichlet series and logarithmic differentiation}

\begin{definition}[Finite Dirichlet convolution]
For arithmetic functions with values in a commutative ring encoded in
naturals or signed pairs, define
\[
 (f*g)(n)=\sum_{d\mid n}f(d)g(n/d).
\]
Let \(\mathbf1(n)=1\), \(\varepsilon(1)=1\) and
\(\varepsilon(n)=0\) for \(n>1\).
\end{definition}

Finite divisor-triple reindexing proves associativity; divisor substitution
proves commutativity; and \(\varepsilon\) is the identity.  Lemma
\ref{lem:mobius-cancel} is exactly \(\mu*\mathbf1=\varepsilon\).

\begin{definition}[Weighted derivation]
For an arithmetic function \(f\), define
\[
 (D_wf)(n)=L_w(n)f(n).
\]
\end{definition}

\begin{theorem}[Leibniz rule]\label{thm:dirichlet-leibniz}
\[
 \boxed{D_w(f*g)=(D_wf)*g+f*(D_wg).}
\]
\end{theorem}

\begin{proof}
For each divisor \(d\mid n\), complete additivity gives
\(L_w(n)=L_w(d)+L_w(n/d)\).  Multiply this equality by
\(f(d)g(n/d)\), sum over divisors, and distribute the finite sum.  The two
resulting sums are precisely the two convolutions on the right.
\end{proof}

\begin{theorem}[Formal logarithmic derivative of zeta]\label{thm:zeta-logder}
In coefficient form,
\[
 \boxed{\Lambda_w=\mu*L_w.}
\]
Equivalently, the formal Dirichlet series of \(\Lambda_w\) is the weighted
logarithmic derivative of the formal zeta series.  For
\(w(p)=\log p\), analytic interpretation in a convergent half-plane gives
\[
 -\frac{\zeta'(s)}{\zeta(s)}
 =\sum_{n\ge1}\frac{\Lambda(n)}{n^s}.
\]
\end{theorem}

\begin{proof}
The coefficient equality is Theorem~\ref{thm:lambda-mobius}.  To see the
logarithmic-derivative language, the zeta coefficient function is
\(\mathbf1\), whose convolution inverse is \(\mu\).  Under the logarithmic
specialization, differentiating \(n^{-s}\) contributes
\(-\log(n)n^{-s}\), so the coefficient function of \(-\zeta'\) is \(L\).
Multiplication by \(1/\zeta\) corresponds to convolution with \(\mu\), giving
\(\mu*L=\Lambda\).  The last sentence is external analytic semantics; the
coefficient identity itself is finite and exact.
\end{proof}

\begin{theorem}[Prime-power coefficients of \(\log\zeta\)]\label{thm:log-zeta-coeff}
There is a rational-coded arithmetic function \(H\) given by
\[
 H(n)=\begin{cases}1/a,&n=p^a\text{ with }p\text{ prime},\\0,&\text{otherwise},\end{cases}
\]
such that \(D_wH=\Lambda_w\).  It is the coefficient function of the formal
Dirichlet logarithm of zeta:
\[
 \log\zeta(s)=\sum_p\sum_{a\ge1}\frac1{a p^{as}}.
\]
\end{theorem}

\begin{proof}
If \(n=p^a\), then \(L_w(n)=aw(p)\), so
\((D_wH)(n)=aw(p)\cdot(1/a)=w(p)=\Lambda_w(n)\); otherwise both sides are
zero.  For a fixed coefficient index \(n\), only exponents
\(a\le\lfloor\log_2 n\rfloor\) can contribute, so the formal logarithm is a
finite rational calculation at that coefficient.  One may derive it from the
finite Euler product over primes \(p\le n\) and the finite truncation
\(-\log(1-X)=\sum_{a=1}^{A}X^a/a+X^{A+1}R_A(X)\), whose remainder has degree
above the target coefficient.  This is status \StatusF.
\end{proof}

\begin{theorem}[Completely multiplicative twist]\label{thm:twist}
If \(f(mn)=f(m)f(n)\) and \(f(1)=1\), then the weighted logarithmic derivative
of the Euler product for \(f\) has coefficient
\[
 \boxed{f(n)\Lambda_w(n).}
\]
\end{theorem}

\begin{proof}
At a prime power \(p^a\), complete multiplicativity gives
\(f(p^a)=f(p)^a\), while formal differentiation of the local geometric factor
selects the coefficient \(w(p)f(p)^a\).  Non-prime-powers cancel under the
formal logarithm.  Equivalently, apply Theorem~\ref{thm:dirichlet-leibniz} to
the convolution inverse of the corresponding coefficient function.  Every
coefficient calculation is over divisors of a fixed \(n\).
\end{proof}

\section{Finite Abel summation and elementary Chebyshev bounds}

\begin{theorem}[Finite Abel summation]\label{thm:abel}
Let \(a(1),\ldots,a(N)\) and \(b(1),\ldots,b(N)\) be coded integer or rational
sequences, and put \(A(t)=\sum_{n\le t}a(n)\).  Then
\[
 \boxed{\displaystyle
 \sum_{n=1}^{N}a(n)b(n)
 =A(N)b(N)+\sum_{n=1}^{N-1}A(n)(b(n)-b(n+1)).}
\]
\end{theorem}

\begin{proof}
Replace \(a(n)\) by \(A(n)-A(n-1)\), distribute the finite sum, shift the
index in the second term, and cancel adjacent terms.  This is a telescoping
identity provable by induction on \(N\).  It is the exact finite core of
partial summation/Mellin-transform arguments involving \(\Lambda\).
\end{proof}

The next theorem records integer inequalities first.  Their logarithmic forms
are obtained with certified rational log bounds.

\begin{theorem}[Central-binomial lcm bounds]\label{thm:binomial-lcm}
For every \(n\ge1\),
\[
 \prod_{\substack{n<p^a\le2n\\p\text{ prime}}}p
 \ \bigm|\ \binom{2n}{n}
 \ \bigm|\ \lcmu(2n).
\]
Moreover,
\[
 \frac{4^n}{2n+1}\le \binom{2n}{n}\le4^n.
\]
All four statements are literal \(\ipra\) theorems.
\end{theorem}

\begin{proof}
Legendre's finite valuation formula gives
\[
 v_p\!\binom{2n}{n}
 =\sum_{j\ge1}
 \left(\left\lfloor\frac{2n}{p^j}\right\rfloor
      -2\left\lfloor\frac n{p^j}\right\rfloor\right),
\]
where the sum stops when \(p^j>2n\).  Each summand is zero or one.  If
\(n<p^a\le2n\), the summand at \(j=a\) is one, proving the left divisibility.
The total valuation is at most the number of powers \(p^j\le2n\), which is
\(v_p(\lcmu(2n))\), proving the right divisibility.

The upper bound follows because \(\binom{2n}{n}\) is one term in
\((1+1)^{2n}=4^n\).  There are \(2n+1\) binomial coefficients whose sum is
\(4^n\), and the central one is maximal: the ratio of successive coefficients
is \((2n-k)/(k+1)\), at least one up to the center and at most one afterward.
Thus the maximum is at least the average, giving the lower bound.  Factorials,
valuations, and binomial coefficients are primitive recursive, and all
arguments are bounded inductions.
\end{proof}

\begin{corollary}[Explicit Chebyshev bounds]\label{cor:chebyshev-bounds}
Under the logarithmic specialization,
\[
 \psi(2n)-\psi(n)\le2n\log2,
\]
and for every integer \(x\ge6\),
\[
 \boxed{\frac{\log2}{3}x\le\psi(x)\le4(\log2)x.}
\]
The native \(\ipra\) content is the integer divisibility and power inequalities
of Theorem~\ref{thm:binomial-lcm}; rational approximations to \(\log2\) give
uniform finite certificates for the displayed real inequalities.
\end{corollary}

\begin{proof}
The left divisibility and \(\binom{2n}{n}\le4^n\) give
\(
 \psi(2n)-\psi(n)\le\log\binom{2n}{n}\le2n\log2.
\)
Choose \(m\) with \(2^{m-1}<x\le2^m\).  Monotonicity and dyadic telescoping
give
\[
 \psi(x)\le\psi(2^m)
 =\sum_{j=0}^{m-1}\bigl(\psi(2^{j+1})-\psi(2^j)\bigr)
 \le\sum_{j=0}^{m-1}2^{j+1}\log2
 <2^{m+1}\log2<4x\log2.
\]
For the lower bound, the right divisibility gives
\(\psi(2n)\ge\log\binom{2n}{n}\).  If \(n\ge3\), then
\(2n+1\le2^n\), so the binomial lower bound yields
\(\psi(2n)\ge n\log2\).  For \(x\ge6\), take
\(n=\lfloor x/2\rfloor\); then \(n\ge3\), \(2n\le x\), and
\(n\ge x/3\).  Monotonicity completes the proof.
\end{proof}

\section{The harmonic Mangoldt sum}

\begin{definition}
For logarithmic weight, let
\[
 A(N)=\sum_{n=1}^{N}\frac{\Lambda(n)}n.
\]
Internally, choose the common denominator \(N!\) and logarithm intervals; the
claim below is represented by inequalities between integer numerators.
\end{definition}

\begin{theorem}[Exact factorial identity]\label{thm:factorial-lambda}
For every \(N\ge1\),
\[
 \boxed{\displaystyle
 \log(N!)=\sum_{d=1}^{N}\Lambda(d)\left\lfloor\frac Nd\right\rfloor.}
\]
More generally, for every weight \(w\),
\[
 L_w(N!)=\sum_{d=1}^{N}\Lambda_w(d)\left\lfloor\frac Nd\right\rfloor.
\]
The weighted identity is literal status \StatusI.
\end{theorem}

\begin{proof}
Sum Theorem~\ref{thm:lambda-divsum} over \(m=1,\ldots,N\):
\[
 \sum_{m\le N}L_w(m)
 =\sum_{m\le N}\sum_{d\mid m}\Lambda_w(d)
 =\sum_{d\le N}\Lambda_w(d)\#\{m\le N:d\mid m\}.
\]
The count is \(\lfloor N/d\rfloor\).  Complete additivity makes the left side
\(L_w(N!)\).  All sums are finite.
\end{proof}

\begin{theorem}[Explicit \(\log N+O(1)\) bound]\label{thm:harmonic-lambda}
For every \(N\ge2\),
\[
 \boxed{\log N-2\log2\le
 \sum_{n\le N}\frac{\Lambda(n)}n
 \le\log N+4\log2.}
\]
Thus the reference-page assertion
\(\sum_{n\le x}\Lambda(n)/n=\log x+O(1)\) has a uniform finite constructive
proof with an explicit constant.
\end{theorem}

\begin{proof}
Since \(\lfloor N/d\rfloor\le N/d\),
Theorem~\ref{thm:factorial-lambda} gives
\[
 \frac{\log(N!)}N\le A(N).
\]
Also \(N/d<\lfloor N/d\rfloor+1\), so
\[
 A(N)<\frac{\log(N!)+\psi(N)}N.
\]
The elementary upper bound \(N!\le N^N\) gives
\(\log(N!)/N\le\log N\), and
The upper-bound half of Corollary~\ref{cor:chebyshev-bounds} is valid for every positive integer \(N\) by the same dyadic argument (and the finitely many values below six can also be checked directly), so
\(\psi(N)/N\le4\log2\).

For the lower factorial estimate, for each \(1\le k\le N\), let \(c(k)\) be
the least \(j\) such that \(N\le2^jk\).  Then
\(k\ge N/2^{c(k)}\), hence
\[
 \log k\ge\log N-c(k)\log2.
\]
For each \(j\ge1\), the number of \(k\) with \(c(k)\ge j\) is at most
\(\lfloor N/2^{j-1}\rfloor\).  Summing these finite level sets gives
\[
 \sum_{k=1}^{N}c(k)
 \le\sum_{j\ge1}\left\lfloor\frac N{2^{j-1}}\right\rfloor<2N.
\]
Therefore
\(
 \log(N!)\ge N\log N-2N\log2
\), and division by \(N\) completes the lower bound.  Internally every
logarithm is replaced by its certified rational interval, and all divisions
are cleared by positive denominators.
\end{proof}

\section{Generalized Mangoldt functions}

One standard generalization is defined by the formal Dirichlet identity
\[
 \frac{(-1)^k\zeta^{(k)}(s)}{\zeta(s)}
 =\sum_{n\ge1}\frac{\Lambda_k(n)}{n^s}.
\]
Its finite coefficient content can be defined without derivatives.

\begin{definition}[Generalized weighted Mangoldt]
For \(k\ge0\), define the signed/rational-coded arithmetic function
\[
 \Lambda_{w,k}=\mu*(L_w)^k,
\]
where \((L_w)^k(n)\) means the ordinary value power, not convolution power.
For logarithmic weight this is the usual generalized Mangoldt function.
\end{definition}

\begin{theorem}
\[
 \Lambda_{w,0}=\varepsilon,
 \qquad
 \Lambda_{w,1}=\Lambda_w,
\]
and every value \(\Lambda_{w,k}(n)\) is obtained by a finite divisor sum
primitive recursive in \((k,n)\) whenever coded exponentiation in the value
ring is available.
\end{theorem}

\begin{proof}
For \(k=0\), the function \((L_w)^0\) is \(\mathbf1\), so
\(\mu*\mathbf1=\varepsilon\).  For \(k=1\),
Theorem~\ref{thm:zeta-logder} gives \(\mu*L_w=\Lambda_w\).  The general value
is by definition
\(\sum_{d\mid n}\mu(d)L_w(n/d)^k\), a bounded sum.
\end{proof}

Further identities involving Riesz means, residues of \(\zeta'/\zeta\), sums
over nontrivial zeta zeros, or equivalences with the Riemann hypothesis are not
asserted as native \(\ipra\) theorems here.  Their finite approximants may be
primitive recursive, but a proof of the claimed limiting or zero-distribution
properties needs a substantial constructive complex-analysis development and,
in the RH cases, would settle an open problem.

\section{Lean formalization architecture}

\subsection{Existing ProveIt infrastructure}

The current \code{ConciseIPRA} development already supplies the pieces needed
for an object-theory formalization:
\begin{itemize}
\item \code{ConciseIPRA.PR}: primitive-recursive symbols \code{PR}, terms
\code{PTm}, the axiom theory \code{PRAAx}, and quantifier-free induction
\code{PRA.qfInd};
\item \code{ConciseIPRA.TeqRec}: a derived primitive-recursion operator
\code{Tm.prc} in the scalar equational theory, with zero and successor
recursion equations;
\item \code{ConciseIPRA.Interpretation}: sound term interpretations and the
preservation theorem \code{Interp.sound};
\item \code{ConciseIPRA.BiInterpretation}: the faithful equivalences
\code{pra\_theorem\_iff} and \code{teq\_theorem\_iff} between ordinary iPRA
and the concise scalar theory.
\end{itemize}
Thus the recommended strategy is to prove the arithmetic theorems first in
\code{PRAAx}, where every primitive-recursive helper is a symbol, and then
transfer the closed formulas to \code{TeqAx}.  Reimplementing all arithmetic
directly in the tiny native calculus would obscure the number theory and
create unnecessary proof-term bulk.

\subsection{Proposed modules}

\begin{longtable}{@{}>{\raggedright\arraybackslash}p{0.34\textwidth}>{\raggedright\arraybackslash}p{0.58\textwidth}@{}}
\toprule
Module & Main contents \\
\midrule
\endfirsthead
\toprule Module & Main contents \\
\midrule
\endhead
\code{IPRA/BoundedFold.lean} & PR compilers for bounded sums, products, counts, minima, list encodings, and their recursion equations.\\
\code{IPRA/Divisibility.lean} & Total division/remainder, divisibility Boolean, gcd/lcm, Bezout trace, congruence, exact quotient.\\
\code{IPRA/Factorization.lean} & Least factor, prime Boolean, valuation, factor trace, reconstruction, uniqueness, prime-power recognizer.\\
\code{IPRA/Mobius.lean} & Signed-pair integers, M\"obius function, squarefree Boolean, cancellation, finite inversion.\\
\code{IPRA/Totient/Basic.lean} & Counting definition, prime-power count, CRT bijection, multiplicativity, product formula.\\
\code{IPRA/Totient/Divisor.lean} & Gcd fibers, divisor sum, M\"obius formula, summatory floor-square identity.\\
\code{IPRA/Totient/Identities.lean} & Power, divisibility, general product, gcd/lcm, doubling, radical, squarefree sum.\\
\code{IPRA/Totient/Units.lean} & Reduced residue permutation, Euler theorem, order, sum of totatives, Menon identity.\\
\code{IPRA/Totient/Applications.lean} & RSA, \(m\mid\varphi(a^m-1)\), bounds, nontotients, finite limsup/liminf witnesses.\\
\code{IPRA/Polynomial.lean} & Integer polynomial codes, division/remainder, cyclotomic recursion, root-of-unity congruence.\\
\code{IPRA/FormalSeries.lean} & Coefficientwise formal power and Dirichlet series, Lambert identity, convolution algebra.\\
\code{IPRA/Mangoldt/Basic.lean} & \(L_w\), \(\Lambda_w\), complete additivity, divisor sum, M\"obius identities.\\
\code{IPRA/Mangoldt/Lcm.lean} & \code{lcmUpto}, valuation, lcm quotient, \(\Psi_w\), factorial identity.\\
\code{IPRA/Mangoldt/Bounds.lean} & Binomial valuations, lcm divisibilities, dyadic bounds, denominator-cleared harmonic estimate.\\
\code{IPRA/Transfer.lean} & Closed theorem formulas and applications of \code{pra\_theorem\_iff}.\\
\bottomrule
\end{longtable}

\subsection{Proof representation discipline}

A useful design is to make every mathematical predicate Boolean-valued at the
PR level.  For example,
\[
 \operatorname{CoprimeB}(a,n)=\ind{\gcd(a,n)=1},
 \qquad
 \operatorname{Phi}(n)=\operatorname{bsum}(n,\lambda a.
 \operatorname{CoprimeB}(a,n)).
\]
The main theorems are then equations between PR terms.  Universal closure of
an equation is quantifier-free, so \code{PRA.qfInd} applies directly.  A
bounded existential such as ``there is a prime-power representation'' is not
left as a logical quantifier; it is compiled into a search function with a
correctness equation.

For finite bijections, avoid cardinality libraries in the object theory.
Encode a map \(f:[0,A)\to[0,B)\) and prove:
\begin{enumerate}
\item a bounded range equation;
\item a Boolean injectivity equation;
\item a bounded inverse equation.
\end{enumerate}
A generic PR theorem then turns these data into equality of bounded counts.
This single lemma supports CRT multiplicativity, gcd fibers, reduced-residue
permutations, and Menon's local count.

\subsection{Illustrative Lean skeleton}

The following is schematic but follows the actual names in the repository:

\begin{lstlisting}
import ConciseIPRA.BiInterpretation
import ConciseIPRA.TeqRec

namespace LeanProofs.ConciseIPRA.NumberTheory
open LeanProofs.ConciseIPRA
open Syn
open scoped Fm

-- Each helper is ultimately a PR symbol assembled by comp/prec.
def phiPR : PR 1 := by
  -- bounded sum of the coprimality characteristic
  sorry

def lambdaWPR (w : PR 1) : PR 1 := by
  -- bounded prime-power search followed by w(base)
  sorry

-- The formula is quantifier-free: equality of two PR applications.
def sumPhiFormula : Fm PTm 1 := by
  -- sum_{d | n} phi(d) = n
  sorry

-- The mathematical induction is performed in PRAAx.
theorem sumPhi_pra : ([] : List (Fm PTm 1)) |- [PRAAx] sumPhiFormula := by
  -- factor trace + gcd-fiber partition + PRA.qfInd
  sorry

-- Faithful transfer to the concise scalar theory.
theorem sumPhi_teq :
    ([] |- [TeqAx] sumPhiFormula.map (fun _ => FromPRA.sigma)) :=
  (pra_theorem_iff sumPhiFormula).mp sumPhi_pra

end LeanProofs.ConciseIPRA.NumberTheory
\end{lstlisting}

The exact parser notation may require the repository's scoped notation to be
opened; the archive's Lean files are deliberately labeled scaffolds rather
than compiled proofs.

\subsection{Host-level validation versus object-theory proof}

Mathlib already contains host theorems such as \code{Nat.totient\_mul},
\code{Nat.sum\_totient}, \code{Nat.totient\_prime\_pow}, and
\code{ArithmeticFunction.vonMangoldt\_sum}.  They are valuable specifications
and regression oracles, but importing them does not prove the result in
\(\ipra\).  The object proof must establish the corresponding encoded formula
from \code{PRAAx}.  The archive separates these roles:
\begin{itemize}
\item \code{VonMangoldtTotientHost.lean} records host-level targets and
cross-checks theorem names;
\item \code{VonMangoldtTotientObjectTheory.lean} gives the staged object-level
obligations and transfer endpoint;
\item \code{FORMALIZATION_PLAN.md} lists every dependency and acceptance
criterion.
\end{itemize}

\section{Constructive audit}

\begin{longtable}{@{}p{0.25\textwidth}p{0.67\textwidth}@{}}
\toprule
Potential hidden principle & Why it is not used \\
\midrule
\endfirsthead
\toprule Potential hidden principle & Why it is not used \\
\midrule
\endhead
Excluded middle & Every case distinction is on a primitive-recursive Boolean: equality, order, divisibility, primality, or bounded membership. Intuitionistic logic proves \(b=0\vee b=1\) for a Boolean computed by recursion.\\
Choice & Prime factors, inverses, orders, and CRT witnesses are selected by explicit least bounded search or Euclidean algorithms.\\
Unbounded minimization & Every minimum has an explicit bound: a divisor is at most the input, an exponent is at most the input, and an order is at most \(\varphi(n)\).\\
Infinite sets & ``Sets'' of divisors, primes, residues, and exponent vectors are bounded characteristic functions.\\
Unique factorization as a black box & A least-factor trace reconstructs the input; valuations prove uniqueness.\\
Rational quotient types & Equalities and inequalities are cross-multiplied using positive denominators; exact natural quotients are accompanied by divisibility proofs.\\
Real logarithm & The exact theory is parameterized by a PR prime weight. Logarithms are finite rational interval specializations at a requested precision.\\
Infinite series rearrangement & Lambert and Dirichlet identities are coefficientwise; a fixed coefficient sees finitely many divisors or multiples.\\
Group cardinality theorems & Finite unit arguments are implemented by explicit permutations, inverses, and bounded counts.\\
Square roots & Bounds are stated internally as squared inequalities or by an integer \(r\) satisfying \(r^2\le n<(r+1)^2\).\\
\bottomrule
\end{longtable}

\section{Statements not claimed}

The reference pages also mention material beyond the present native scope.
This report does not claim:
\begin{itemize}
\item the prime number theorem, explicit formulas summing over nontrivial zeros
of \(\zeta\), RH-equivalent error terms, or oscillation theorems for Riesz
means;
\item Carmichael's conjecture on the multiplicity of totient values;
\item the full classification and density theory of totients/nontotients;
\item asymptotic constants involving \(\pi\), \(e^\gamma\), or infinite Euler
products unless replaced by the exact finite identities proved above;
\item a completed, kernel-checked Lean formalization of the object theory.
\end{itemize}
These exclusions prevent analytic notation or a host theorem from being
misreported as an \(\ipra\) derivation.

\section{Conclusion}

The arithmetic theories of \(\varphi\) and \(\Lambda\) are exceptionally well
suited to \(\ipra\).  Their core is finite factorization, bounded counting,
and finite convolution.  Once logarithm is abstracted to a prime weight, the
Mangoldt formulas become exact primitive-recursive identities.  Once unit
groups are replaced by explicit reduced-residue permutations, the central
totient theorems become quantifier-free induction arguments.  Formal power
series and Dirichlet series add no infinitary burden when interpreted
coefficientwise.

The remaining work is therefore engineering rather than a missing
foundational idea: build a reusable bounded-fold and factorization library in
\code{PRAAx}, prove the finite bijection and convolution lemmas once, instantiate
them for the two functions, and invoke the existing faithful
\code{PRAAx}\(\leftrightarrow\)\code{TeqAx} transfer.  The module plan in this
archive gives a direct path to a complete Lean development.

\begin{thebibliography}{99}
\bibitem{wiki-mangoldt}
Wikipedia contributors, ``Von Mangoldt function,'' accessed 15 November 2026,
\url{https://en.wikipedia.org/wiki/Von_Mangoldt_function}.

\bibitem{mw-mangoldt}
Eric W. Weisstein, ``Mangoldt Function,'' MathWorld, accessed 15 November 2026,
\url{https://mathworld.wolfram.com/MangoldtFunction.html}.

\bibitem{wiki-totient}
Wikipedia contributors, ``Euler's totient function,'' accessed 15 November 2026,
\url{https://en.wikipedia.org/wiki/Euler%27s_totient_function}.

\bibitem{mw-totient}
Eric W. Weisstein, ``Totient Function,'' MathWorld, accessed 15 November 2026,
\url{https://mathworld.wolfram.com/TotientFunction.html}.

\bibitem{proveit}
Vladimir Reshetnikov, \emph{ProveIt}, revision
\code{5fb2ca5dda188a4a3c7ce9ba6f28030bd0fe8322}, 15 November 2026,
\url{https://github.com/VladimirReshetnikov/ProveIt}.

\bibitem{mathlib-totient}
Mathlib contributors, \code{Mathlib/Data/Nat/Totient.lean}, inspected at
revision \code{cbe1882f5b62496bdfc6bbe540a2b3f2db852440}.

\bibitem{mathlib-mangoldt}
Mathlib contributors,
\code{Mathlib/NumberTheory/ArithmeticFunction/VonMangoldt.lean}, inspected at
the same revision.

\bibitem{mathlib-chebyshev}
Mathlib contributors, \code{Mathlib/NumberTheory/Chebyshev.lean}, inspected at
the same revision.
\end{thebibliography}

\appendix
\section{Compact theorem checklist}

\begin{longtable}{@{}>{\raggedleft\arraybackslash}p{0.04\textwidth}>{\raggedright\arraybackslash}p{0.68\textwidth}>{\raggedright\arraybackslash}p{0.18\textwidth}@{}}
\toprule
 & Claim & Location/status \\
\midrule
\endfirsthead
\toprule & Claim & Location/status \\
\midrule
\endhead
1 & \(\varphi(n)=\#\{a<n:\gcd(a,n)=1\}\) & Def.~\ref{def:totient}, \StatusI\\
2 & \(\varphi(p)=p-1\), \(\varphi(p^a)=p^{a-1}(p-1)\) & Thm.~\ref{thm:phi-prime-power}, \StatusI\\
3 & Coprime multiplicativity and product formula & Thm.~\ref{thm:phi-product}, \StatusI\\
4 & Gcd-fiber count & Thm.~\ref{thm:gcd-fiber}, \StatusI\\
5 & \(\sum_{d\mid n}\varphi(d)=n\) & Thm.~\ref{thm:sum-phi}, \StatusI\\
6 & M\"obius formulas for \(\varphi\) & Cor.~\ref{cor:phi-mobius}, \StatusI\\
7 & Summatory floor-square identity & Thm.~\ref{thm:summatory-phi}, \StatusI\\
8 & Farey count & after Thm.~\ref{thm:summatory-phi}, \StatusI\\
9 & \(\varphi(n^k)=n^{k-1}\varphi(n)\) & Thm.~\ref{thm:phi-power}, \StatusI\\
10 & \(a\mid b\Rightarrow\varphi(a)\mid\varphi(b)\) & Thm.~\ref{thm:phi-divides}, \StatusI\\
11 & General multiplication formula & Thm.~\ref{thm:phi-general-mul}, \StatusI\\
12 & Gcd-lcm identity & Cor.~\ref{cor:phi-gcd-lcm}, \StatusI\\
13 & Doubling formula & Thm.~\ref{thm:phi-double}, \StatusI\\
14 & Radical identity & Thm.~\ref{thm:phi-rad}, \StatusI\\
15 & Squarefree-divisor identity & Thm.~\ref{thm:squarefree-divisor}, \StatusI\\
16 & Evenness and \(2^r\mid\varphi(n)\) & Thm.~\ref{thm:phi-even}, \StatusI\\
17 & Sum of totatives & Thm.~\ref{thm:sum-totatives}, \StatusI\\
18 & Menon's identity & Thm.~\ref{thm:menon}, \StatusI\\
19 & Euler/Fermat theorem & Thm.~\ref{thm:euler}, \StatusI\\
20 & Order divides totient & Cor.~\ref{cor:order-divides}, \StatusI\\
21 & \(m\mid\varphi(a^m-1)\) & Thm.~\ref{thm:order-am1}, \StatusI\\
22 & RSA correctness & Thm.~\ref{thm:rsa}, \StatusI\\
23 & Cyclotomic gcd-sum identity & Thm.~\ref{thm:cyclotomic-gcd-sum}, \StatusF\\
24 & Lambert series & Thm.~\ref{thm:lambert}, \StatusF\\
25 & Totient Dirichlet series coefficients & Thm.~\ref{thm:phi-dirichlet}, \StatusF\\
26 & Lower and upper bounds & Thms.~\ref{thm:phi-lower}, \ref{thm:phi-upper-composite}, \StatusI\\
27 & Totient-sigma identity and \(1/2\) bound & Thm.~\ref{thm:phi-sigma}, \StatusI\\
28 & Finite limsup/liminf witnesses & Thms.~\ref{thm:limsup-phi}, \ref{thm:liminf-phi}, \StatusI/\StatusA\\
29 & Odd nontotients and \(\varphi(n)=2\) & Thms.~\ref{thm:odd-nontotient}, \ref{thm:phi-two}, \StatusI\\
30 & Weighted Mangoldt definition & Def.~\ref{def:weighted-lambda}, \StatusI\\
31 & Mangoldt divisor sum & Thm.~\ref{thm:lambda-divsum}, \StatusI\\
32 & M\"obius forms & Thm.~\ref{thm:lambda-mobius}, \StatusI\\
33 & M\"obius--Mangoldt convolution & Thm.~\ref{thm:mu-lambda}, \StatusI\\
34 & Lcm quotient & Thm.~\ref{thm:lcm-quotient}, \StatusI/\StatusA\\
35 & \(\Psi_w=L_w(\lcmu)\) & Thm.~\ref{thm:psi-lcm}, \StatusI\\
36 & Dirichlet Leibniz rule & Thm.~\ref{thm:dirichlet-leibniz}, \StatusF\\
37 & \(-\zeta'/\zeta\) coefficient identity & Thm.~\ref{thm:zeta-logder}, \StatusF/\StatusA\\
38 & Coefficients of \(\log\zeta\) & Thm.~\ref{thm:log-zeta-coeff}, \StatusF\\
39 & Finite Abel summation & Thm.~\ref{thm:abel}, \StatusI\\
40 & Binomial-lcm divisibilities and Chebyshev bounds & Thm.~\ref{thm:binomial-lcm}, Cor.~\ref{cor:chebyshev-bounds}, \StatusI/\StatusA\\
41 & Factorial Mangoldt identity & Thm.~\ref{thm:factorial-lambda}, \StatusI\\
42 & \(\sum\Lambda(n)/n=\log N+O(1)\) with constants & Thm.~\ref{thm:harmonic-lambda}, \StatusI/\StatusA\\
43 & Generalized \(\Lambda_k\) and \(\Lambda_1=\Lambda\) & Generalized Mangoldt section, \StatusI/\StatusF\\
44 & Subbarao congruence classification & Thm.~\ref{thm:subbarao}, \StatusI\\
45 & \(n\mid\varphi(a^n+b^n)\) & Thm.~\ref{thm:plus-power}, \StatusI\\
46 & Goldbach-to-totient implication & Prop.~\ref{prop:goldbach-phi}, conditional \StatusI\\
\bottomrule
\end{longtable}

\section{Repository and reproducibility notes}

The archive was prepared against the ProveIt main-branch revision recorded in
the bibliography.  The current \code{ConciseIPRA} README describes the
calculus as bi-interpretable with full first-order iPRA; \code{PR.lean}
contains the ordinary iPRA syntax and quantifier-free induction;
\code{TeqRec.lean} derives primitive recursion in the scalar theory; and
\code{BiInterpretation.lean} proves faithful theorem transfer.  The host
mathlib files were used only as theorem inventories and cross-checks.

The finite Python checker is not a proof of the universal theorems.  It is a
regression artifact that tests definitions, identities, examples, and the
polynomial congruence over a bounded range.  The universal proofs are the
arguments above; a complete formal certification requires implementing the
Lean roadmap.

\end{document}
